\documentclass[../../../include/open-logic-section]{subfiles}

\begin{document}

\olfileid{his}{set}{mythology}

\olsection{Luwih Mitos tinimbang Sajarah?}

Yèn nyawang manèh prastawa-prastawa iki
saka jarak luwih saka satus taun,
kita kudu ngati-ati supaya ora
langsung nampa putusan-putusan
mau tanpa mriksa. Asil-asilé
pancen anyar, nyenengaké, lan
nggumunaké. Nanging sepira
gedhéné geger kang ditimbulaké?
Lan apa asil-asil iku tenan
mbuktèkaké yèn kita ora
kena ngandelaké intuisi géomètri?

Ngenani geger mau, \citet{Gouvea2011}
nuduhaké yèn cathetan Cantor kang
kondhang marang Dedekind,
``\emph{je le vois, mais je ne le crois
pas}'', kerep dipetik tanpa
konteks kang cukup. Iki
pethikan kang luwih jangkep
(dikutip saka \citeauthor{Gouvea2011}):
\begin{quote}
Nyuwun pangapunten marang semangatku
ngenani prakara iki yèn aku kerep
njaluk kabecikan lan kasabaranmu;
pawarta kang lagi waé takkirimi
marang kowé uga ora taksangka
lan anyar banget kanggo aku,
saéngga aku ora bisa ayem
nganti nampa putusanmu,
kancaku kang takajèni,
ngenani kabenerané.
Sadurungé kowé sarujuk
karo aku, aku mung
bisa kandha:
\emph{je le vois, mais je ne le crois pas.}
\end{quote}
Cantor ngerti yèn asilé ``ora
disangka-sangka lan anyar banget''.
Nanging ora cetha yèn dhèwèké
nganggep asil iku \emph{ora bisa
dipercaya}. Kaya kang diterangaké
\citeauthor{Gouvea2011}, dhèwèké
mung njaluk Dedekind mriksa
bukti kang wis diajokaké.

Ngenani intuisi géomètri: Peano
nerbitaké kurvané kang ngebaki
ruang tanpa gambar. Nanging
nalika Hilbert nerbitaké kurvané,
dhèwèké nerangaké tujuwané:
dhèwèké arep menehi para
maca cara kang cetha kanggo
mangertèni asil Peano,
yèn para maca ``migunakaké
intuisi géomètri ing ngisor
iki''; banjur dhèwèké
nglebokaké rerangkèn
\emph{gambar} kaya kang
ana ing
\olref[his][set][pathology]{sec}.

Luwih umum manèh: senajan
gambar rada arang katon
ing bukti kang diterbitaké,
ora bisa disangkal yèn
para matématikawan
\emph{kerep} nganggo
gambar nalika mbuktèkaké
prakara. (Kanthi kasar:
matématikawan kang apik
ngerti kapan intuisi
géomètri bisa diandelaké.)

Mula aja gampang percaya marang
gegap-gempitané; paling ora,
aja mung nampa tanpa mriksa.
Kanggo katrangan luwih,
delengen \citet{Giaquinto2007}.

\end{document}
